% !TeX root = ../papers/nc-p3.tex
\section{滑らかな種数一核曲線を持つ八余因子窓の排除}
\label{sec:smooth-elliptic-cofactor-exclusion}
\begin{theorem}\label{thm:smooth-elliptic-cofactor-exclusion}
四辺有限な標準有限Frobenius窓で、整な反対八次曲面の余因子の包が八次元なら、
核曲線は必ず特異である。節点または尖点のある整な二次式完全交叉に限られる。
この排除は点グラフの各局所長に条件を課さない。
\end{theorem}
\begin{proof}[分岐元における二重数上の核]
滑らかな場合を仮定し、前節の二捩れの消滅元をまず扱う。
$P=L\eta$、$\eta\ne1$、$P^2=L^2$ としてその対称評価行列を $C_0$ とする。
一般の $x$ で階数1以下ではない。$H^0(P)$ の平方における階数1の点で
二次関係二本を満たすものは、二次式完全交叉である $E$ の評価そのものなので、
$L$ での対合も階数1になってしまう。
階数2なら一般の左核は $H^0(P)$ の鉛筆で、前節の固定因子の議論により基点を持たない。
その二基底 $a,b$ の積による零化直線の逆は前節の $\mathcal B(a,b)$ と同じ公式である。
今は $\eta=\eta^{-1}$ なので $K(w,z)=-K(z,w)$ である（留数と極を消した切断の消滅）。
従って $\mathcal B$ は $a,b$ に関し交代であり、元の $N(x)$ の値は全て
$\wedge^2 H^0(P)$ の固定した六次元以下の像に入る。十成分の独立性に反する。
よって $C_0$ の一般階数は3となる。

$P_t$ を $P$ でのPicard曲線の形式パラメータによる変形とし、
$P_tP_{-t}=L^2$ とする。反転で符号の変わる局所パラメータと互いに転置される切断基底を取る。
代入した対合は $C_t=C_0+tC_1+\cdots$、$C_1$ は交代であり、
$\det C_t$ は $t^2$ で割れる。$D$ の一般点での一点評価は全ての $t$ で
階数1なので $K_t=\operatorname{adj}C_t/d^2$ は $t$ の各係数で二次式行列となる。
これは $d$ の一般点で単元一成分を消去するか、全切断束上の多項式除法で従う。
$K_0=aa^{\mathsf t}$ と線型列で書ける。一定の核なら対応する切断と全切断との積を
八座標全てが消すことになり、原始分解で定数列は不可だからである。
この線型写像の階数は3以上である。もし2なら $[a]$ の一般逆像は入力平面で、
その平面上、余因子汎関数は固定された四点因子 $D_a$ 上に支えられる。
それらの可逆な $L$ 対合の逆の包は高々四次元である。実際 $L$ の評価は $D_a$ 上同型で、
一つの可逆対合を固定すれば他は四点環内の単元倍、その逆も逆単元倍となる。
一方、平面上の $N$ の値は全二次式から来て六次元を張るため矛盾する。

$K'_0$ は交代で $C_0K'_0=-C_1aa^{\mathsf t}$ より、$a$ の直交超平面を
$a$ の直線に送る。従って分数体上で $a\wedge K'_0=0$。
$a$ の成分には独立な一次式が三つ以上あるので、Koszul完全性により
$K'_0=ba^{\mathsf t}-ab^{\mathsf t}$ となる一次式列 $b$ がある。
三つだけ独立なら零の第四成分を分け、三元Koszul複体の外次数1,2の完全性を使う。
これで $C_0b+C_1a=0$、すなわち $a_t=a+tb$ は右核の一次持上げである。

二重数に制限した $P_t$ をもとの $E$ 上の階数2束 $F$ と読むと、
$0\to P\to F\to P\to0$ は非分裂である。遷移関数の一次係数が
Picard曲線の非零接方向を表すためである。$a_t$ は $x$ に線型な切断 $s(x)$ を与える。
余因子汎関数は $\det(s,H^0(F))\subset H^0(P^2)$ を消す。
実際任意の持上げ $h_t=h+tk$ と反対向きの $h_{-t}$ に対し
$a_t h_{-t}$ の一次部は局所自明化で $bh-ak$、すなわち行列式の符号倍である。
一般の $x$ で $s$ は零点を持たない。零因子を $A$ とすると、その飽和直線束
$\mathcal O(A)\subset F$ から、$A>0$ なら第一コホモロジー消滅で
行列式の像は $A$ 上で消える全切断となる。さらに $\deg A<4$ である。
$P$ への射が非零で次数4なら同型となり拡張が分裂するからである。
すると $L$ 対合の階数は $\deg A<4$ となり矛盾する。零点がなければ
$0\to\mathcal O\xrightarrow{s}F\to L^2\to0$ から行列式の像は超平面で、
これが余因子汎関数を一意に指定する。
\end{proof}

\begin{proof}[分岐元の二次式表示と加法群トーサー]
$s$ の二次式行列を
\[
 B_F(s)=\left.\frac{\mathcal B_t(a_t,a_{-t})}{t}\right|_{t=0}
\]
と定める。ここで $\mathcal B_t$ には $P_t,P_{-t}$ の核を使い、$a_t$ の一次までの
切断だけが必要である。核は $\eta\ne1$ の近傍で正則に選べる。
定数項は交代性により零である。
一般の $a$ は単純四根を持つ。階数3以上の線形系には固定二重点がなく、
固定部を除いた一般切断は標数零で単純である。
動く四根での前節の対角公式を微分すると、$B_F(s)$ の四対角値は
$a'(z)$ と $s$ の $D_a$ 上の残る成分との積の共通非零定数倍になる。
実際 $a_t$ の根で $a_{-t}$ を評価した一次係数は、局所自明化での
$s=(a,b_{\rm loc})$ の $-2b_{\rm loc}$ である。
従って一般の零点なし $s$ でこれは可逆で、逆の評価汎関数は
$\det(s,H^0(F))$ を消す。後者は例えば前節の二つの積の零化を、和と一次差に取って
極限すれば従う。よって $B_F(s(x))$ は $N(x)$ に定数比例する。
ともに二次式で後者の成分は原始だからである。$a(x)=0$ ではこの微分も零なので、
$a$ は同型である。$s$ の選択は $H^0(F)\to H^0(P)$ の線型切断として読める。

重根と非被約な底にも使うため、$B_F$ の行列式は $s$ の零点の結果式であることを確認する。
$s(z)=0$ なら微分核に $z$ を代入して零となる。零点を印した全切断の入射は射影束で既約、
その像因子の次数は8である。実際一般鉛筆の二切断の外積は $\det F$ の八零点を持ち、
各点でただ一つの鉛筆切断が消える。一般には単純で、同じ切断に二零点はない。
これは $F$ の任意の長さ2因子上の評価が全射であることから、各二重零や二点同時消滅の
入射の次元を数えても分かる（次数2を引いた $P$ の第一コホモロジーは零）。
非零な八次行列式はこの因子と等しい。$a\ne0$ で $D_a$ 上に零のない参照持上げ
$s_0$ を取り、$L$ の $D_a$ への評価で同定する。
$B_F(s)$ は $B_F(s_0)$ に、$D_a$ 上の残る成分の比の乗法を出力側で合成したものである。
単純四根の微分対角公式の等号が普遍平坦因子上で延びるためである。
よって $X$ 上で $s$ の零点は平坦長さ3の共通因子 $A$ となり、$D_a=A+u$ と残余を取れる。

$u$ を固定し $A_u=H^0(P(-u))$ とする。$P_t/P$ に対応する非零の一次移動を $t v$ と書く。
$s$ が因子 $A$ で消えることは、$a_t$ が定数変形した $A$ で消えることに等しい。
残余点の一次移動は $u_t=u+tv$ に限られる。次数1線束のAbel写像が同型だからである。
$P_t(-u_t)$ は $P(-u)$ の定数変形と各ファイバーで同型である。
これを $u$ 全体で比較すると底からの一次変形線束 $J_t$ が差となり、
$J_0=\mathcal O$、その一次類は非零である。実際曲線変数の原点で普遍因子の比を
制限すれば $\mathcal O([o]-[o-tv])$ を得る。
従って特殊化が恒等な二つの線束の同定は、底 $E$ 上の非自明な加法群トーサーをなす。
同定を一つ印すと、$A_u$ の元から $P_t$ の一次持上げへの線型写像 $j_{u,\lambda}$ が決まる。
定数変形した $P(-u)$ の切断を同定で移し、$u_t$ で消える包含を用いるためである。
同定を $\lambda$ だけ変えると、その微分は $A_u\hookrightarrow H^0(P)$ の $\lambda$ 倍だけ変わる。
$X$ 上では持上げの二切断が同じ相対因子 $A$ を持つので、特殊化が1の比例単元で一致し、
この同定をスキーム上で一意に復元する。

トーサーを $\Sigma=\mathbb P(\mathcal E)$ と完備化する。
$\mathcal E$ は自明線束同士の拡張であり、無限遠で
$s|_{A_u}-j_{u,\lambda}$ の微分差の先頭項は包含 $A_u\hookrightarrow H^0(P)$ である。
従ってこれは $A\otimes\mathcal O_\Sigma(-1)\to H^0(P)\otimes\mathcal O_\Sigma$
という四行三列写像で、その退化は必ず有限である。
垂直曲線は無限遠を避けられず、水平曲線があればその正規化への有限引戻しで
$J_t$ の一次類が零となるが、跡を掛け戻すと非零類の写像次数倍なので矛盾する。
前節の次数2の複体の余核長はここでも16である。具体的には $\deg A^*=4$、
$\deg\operatorname{Sym}^2 A^*=16$、$\pi_*\mathcal O(1),\pi_*\mathcal O(2)$ は
階数2,3で次数零であり、楕円曲線上のEuler数で $3\cdot16-4\cdot2\cdot4=16$ となる。
階数0は入力平面を整な $D$ に含ませるので排除される。
従って前節のファイバー上の二次切断と中山の補題が $X$ にも適用でき、長さ20に反する。
これで分岐元は全て除かれた。
\end{proof}

\begin{proof}[四捩れの水平二重被覆と垂直部分]
前節の非分岐表示で $\delta$ が位数2の場合を扱う。
スカラーの底 $\Sigma$ にある水平曲線は、二乗が一定の連結エタール二重被覆
$\widetilde E\to E$ である。この曲線は一つであり、一般点で重複度1である。
一般のスカラーファイバーの退化は高々三次だからである。一般点の核は直線である。
核が二次元なら四次ではなく三次の最大行列式が各スカラー根で二重以上に消え、
共役な二根を合わせると次数を超える。

$A_u$ の $T$ による像から $H^0(Q)/A'_u$ への写像の核束を $K$ とする。
この写像は恒等零ではない。全 $u$ で超平面が一致すれば評価の楕円四次曲線が
$u\mapsto u-\delta$ により射影同値となり、$P\simeq t_{-\delta}^*Q$ が必要である。
右辺は $4\delta=0$ より $Q$ と同型であり $P\ne Q$ に反する。
従って $K$ は飽和階数2部分束、$A/K$ は線束である。この商から
$H^0(Q)/A'_u$ への写像の零因子を $S$、その次数を $k$ とする。
$\deg A=-4$、$\deg(H^0(Q)/A')=4$ より $\deg K=-8+k$ である。

$T|_K$ は $A'=A\otimes J$ の同定で $K\otimes J$ へ写る。
一般点で水平被覆の二固有直線が $K$ を張り、その上でこの主張が成立し、
$K$ の飽和性により全体へ延びる。得る $K\to K\otimes J$ は、$J^2$ の自明化で
平方が定数倍の恒等なので同型である。被覆上では $K$ は二つの相異なる固有線束に分かれ、
互いに被覆変換で移る。各固有直線は入力点を与え、一方で定めた
$\widetilde E\to\mathbb P^3$ の超平面次数は従って
\[
                      b=-\deg K=8-k .
\]
この射は非定数である。一つの $a$ が全ての $u$ で消えることはない。
$\mathbb P(K)$ 内のスカラー付き入射は、この二固有直線の被覆そのもので滑らかである。

局所の残余長を数える。底の完備局所環で $J$ を自明化し、$K$ の二固有値を
$\lambda_+,\lambda_-$ とする。差は単元である。
二固有ベクトルと一補列で $A$ 上の直線の座標を $(p,q,r)$ とすれば、全入射の方程式は
\[
 (\lambda-\lambda_+)p=f_+r,\qquad(\lambda-\lambda_-)q=f_-r,\qquad
 (\lambda-h)r=0,\qquad sr=0 .
\]
ここで $s$ は $S$ の局所方程式であり、$f_+,f_-,h$ は補列の $T$ による成分である。
各 $u\in S$ で $n_u=\operatorname{ord}_u(s)$ とする。
$h$ が両固有値と異なれば $e_u=0$ とする。一方と一致する場合を $+$ と書き
\[
        e_u=\min(n_u,\operatorname{ord}_u(h-\lambda_+),
                        \operatorname{ord}_u(f_+))
\]
と置く。一致はこの点での値を指し、恒等零の位数は無限大とする。

$X$ 上の $r$ の生成するイデアルの長さは $k+2\sum e_u$ 以下である。
このイデアルは局所では $\mathcal O_X/\operatorname{ann}r$ に同型で、
$s=0,\lambda=h$ の上のスキームになる。$R=\mathcal O_{E,u}/(s)$ 上で
$q$ を消去すると、残る線型座標加群は
$R\oplus R/(t^{e_u})$ となる。$t$ は局所パラメータであり、
二変数に課す一行 $(h-\lambda_+,-f_+)$ の単因子を取った。
両固有値と異なるときは自由な一座標だけで同じ式である。
この加群の対称二乗の長さは $n_u+2e_u$ である。ファイバーは入力の点または直線上にあり、
直線切断長は3以下なので、二次式からファイバーの全切断へ全射となる。
中山の補題で厚みを含む全長の上界が従う。

一方 $r=0$ の閉部分は先の $\widetilde E$ に含まれる。
三次式の $X$ 上の像の次元は、その被覆上での次元と $r$ イデアルの長さの和以下なので
\[
                         20\ \le\ 3b+k+2\sum e_u .                 \tag{*}
\]
ここで $3b>0$ なので楕円曲線上の三次制限線束の切断次元が $3b$ である。
\end{proof}

\begin{proof}[一般の母線での三重軌跡の長さ]
一般の $z\in E$ で $a(z)=b'(z)=0$（この式の $b'=Ta$ は第二区分の切断）という
入力直線 $\ell_z$ を取る。二条件が比例する場合は入力平面が $D$ に入ってしまう。
$\ell_z$ 上で $N$ は定数の評価ベクトル $z$ を消すので、残りの対称三行二次式行列を取る。
その六次行列式は非零である。実際 $D$ は一般的に階数3であり、結果式の共通点を印した
母線族が $D$ を覆うので、一般の母線が丸ごと階数2以下に入ることはない。
従って $\ell_z$ 上の階数1以下の軌跡の長さは高々3である。
直線の各局所離散付値環で最小単因子は単元（全二次成分は共通基点を持たない）、
次の単因子の位数が階数1軌跡の長さであり、六次行列式の位数はその二倍以上である。

$\widetilde E$ 上で $a(z)$ の零因子は次数 $b$ で、そのうち $u=z$ のファイバーは次数2である。
残る共通三因子に属する部分は次数 $b-2$ で、$u$ と $v$ は常に異なるので
$\ell_z$ との交叉にこの長さが入る。これは残余因子の族の分解を
$\{z\}\times\widetilde E$ に制限した式であり、重複度も含む。

さらに各 $e_u>0$ に一つの垂直直線があり、一般の $\ell_z$ への寄与は $e_u$ である。
実際 $r\ne0$ のチャートで $\lambda=h$、$q/r$ は正則に消去でき、
$w=p/r$ には $s=0,(h-\lambda_+)w-f_+=0$ だけが残る。
一般の $z$ では $a(z)=0$ が $w$ を一意に消す。底の閉点の固定された
入力直線上の $a$ の二基底は独立な楕円切断なので、この $w$ の値は $z$ とともに非定数に動く。
従って最初の非零係数が相殺される例外値を避け、商の長さはちょうど $e_u$ である。
$z\ne u,v$ なので共通三因子の二表示から $b'(z)=0$ もこの商上で成立する。
前節の乗法による階数1判定と残余の一意な復元は相対スキームでの同型なので、
これらの長さは入力直線でも失われない。
一般 $z$ で垂直の寄与は水平曲線と交わる有限な端点を避け、互いにも重ならない。
従って
\[
                    b-2+\sum e_u\ \le\ 3 .
\]
これと $b+k=8$ を $(*)$ に入れると右辺は高々18となって矛盾する。
四捩れも不可能であり、前節の全候補を尽くした。
\end{proof}

\path{computations/verify_elliptic_branch_member.py} は $y^2=x^3-x$ の二捩れ $(0,0)$ で
積写像の一次係数を直接計算する。一般の線型持上げの有理標本について、
十個の二次成分と八個の六次余因子を持つ微分Bezout行列を良い素数で復元する。
この標本の二次小行列式は十九独立であり、有限辺の四次評価と両立しない。
